\scanpage{36}

\uline{Remarque 3} Les conditions $LG$ et $LG$ effective sont
stables par passage de $X$ à un voisinage ouvert $U$ de $Y$.
Il n'en est plus de même en général des conditions $(LG)_n$,
\deleted{en particulier $(LG)_0$ \unsure{Cependant}
on peut prouver facilement qu'il en est encore ainsi de
$(LG)_0$, [et d'ailleurs]}
probablement pas même pour $(LG)_0$, en raison du fait
que un module loc. libre \inserted{$F$} sur un voisinage $U$
de $Y$ n'est pas restriction d'un module loc libre sur
\deleted{un voisinage} $X$.
Si cependant $\forall$ tel $F$ est isom. à un \uline{quotient}
d'un module loc. libre sur $X$ [p. ex. si $\forall$ Module
coh. sur $X$ est quotient d'un Module loc libre,
p. ex. $X$ quasi-projectif]]
\uline{alors} $(LG)_0$ pour $(X,Y)$ implique
$(LG)_{\unsure{n|0}}$ pour $\forall$ $(U,Y)$ i.e. $LG$ pour $(X,Y)$.
En effet on est ramené à prouver que pour $F$ loc libre
\uline{sur $X$},
\[
H^0(\unsure{U},F)\longrightarrow H^0(U_{/Y},F_{/Y})
\]
est bijectif, on sait déjà qu'il est surjectif,
[car
\diagram{\begin{tikzcd}[column sep=small]
H^0(X,F)\arrow[r,dash] & H^0(U,F)\arrow[r,dash] & H^0(U_{/Y},F_{/Y})
\end{tikzcd}}
l'est] reste à vérifier l'\uline{injectivité}, ce qui résulte de
\deleted{\unsure{pour}} \deleted{\unsure{donc que}}

\margin{35}
